%% kapa: exam=1 round=32 year=2021 period=2 form=A subject=appraisal_law subject_ko=감정평가관계법규 seq=21 number=21
%% source: Q-Net 감정평가사 기출문제 (cst00309) artlSeq=5208397 file=extracted/1008/5208397/A형/2021년도 제32회 감정평가사 1차 2교시 A형.pdf
%% conversion: kordoc 4.24.51 -> tools/gen1.py (auto); review=proofread-sonnet; flags=none
\begin{question}{21}
국유재산법령상 행정재산에 관한 설명으로 옳지 않은 것은?
\begin{choices}
  \item 중앙관서의 장은 행정재산을 직접 공공용으로 사용하려는 지방자치단체에 사용허가하는 경우에는 사용료를 면제할 수 있다.
  \item 중앙관서의 장은 사용허가를 받은 행정재산을 천재지변으로 사용하지 못하게 되면 그 사용하지 못한 기간에 대한 사용료를 면제할 수 있다.
  \item 중앙관서의 장은 행정재산의 사용허가를 철회하려는 경우에는「행정절차법」 제27조의 의견제출을 거쳐야 한다.
  \item 중앙관서의 장은 행정재산으로 사용하기로 결정한 날부터 5년이 지난 날까지 행정재산으로 사용되지 아니한 경우에는 지체 없이 그 용도를 폐지하여야 한다.
  \item 행정재산은「민법」 제245조에도 불구하고 시효취득의 대상이 되지 아니한다.
\end{choices}
\end{question}
\answer{3}
